\section{\S 3.1 水的电离和溶液酸碱性}\label{sect:水的电离和溶液酸碱性}

    \subsection{水的电离}\label{subsec:水的电离}
        
        \cemh{H2O + H2O <=> H2O+ + OH^-}，
        简写为 \cemh{H2O <=> H+ + OH^-}

        \textcolor{red}{往纯水里加入酸或碱，水的电离平衡如何移动？}

        \textcolor{red}{\(c(\cemh{H+})\)与\(c(\cemh{OH^-})\)如何变化？}

        水溶液中\(c(\cemh{H+})\)与\(c(\cemh{OH^-})\)\textcolor{red}{此消彼长}

        一定温度下，纯水或稀溶液中：\(K_\text{W} = \cemh{[H+]} \cdot \cemh{[OH^-]}\)
        \footnote{\cemh{[H2O]} 在稀溶液/纯水中近似为常数}

        \(K_\text{W}\)——水的离子积常数\qquad 简称：水的离子积

        \SI{25}{\degreeCelsius}时： \( K_\text{W} = \SI{1E-14}{}\)\quad\qquad \textcolor{red}{\( K_\text{W} \) 随温度如何变化？}

        \SI{99}{\degreeCelsius}时： \( K_\text{W} = \SI{1E-12}{}\)\quad\qquad \textcolor{blue}{温度变化不大时，使用\SI{25}{\degreeCelsius}时的\( K_\text{W} \)}

    \subsection{溶液酸碱性与\texorpdfstring{\cemh{[H+]}}{H⁺}、\texorpdfstring{\cemh{[OH^-]}}{OH⁻}及 pH 的关系}\label{subsec:溶液酸碱性与[H+]、[OH^-]及pH的关系}

        \( \cemh{pH} = -\lg{}\cemh{[H+]} \) 

        \noindent
        \begin{tabular}{ccccc}
                        & 任何条件下 & & 常温下（ \( K_\text{W} = \SI{1E-14}{}\) ）\\
            中性溶液 & \( \cemh{[H+]} = \cemh{[OH^-]} \) & \( \cemh{pH} = 7 \) & \textcolor{blue}{\( \cemh{[H+]} = \cemh{[OH^-]} = \SI{1E-7}{\mole\per\litre} \) } \\
            酸性溶液 & \( \cemh{[H+]} > \cemh{[OH^-]} \) & \( \cemh{pH} < 7 \) & \textcolor{blue}{\( \cemh{[H+]} > \SI{1E-7}{\mole\per\litre} > \cemh{[OH^-]} \) } \\
            碱性溶液 & \( \cemh{[H+]} < \cemh{[OH^-]} \) & \( \cemh{pH} > 7 \) & \textcolor{blue}{\( \cemh{[H+]} < \SI{1E-7}{\mole\per\litre} < \cemh{[OH^-]} \) }
        \end{tabular}

        \textcolor{blue}{\cemh{[H+]} 越大，溶液酸性越强；\cemh{[H+]} 越小，溶液酸性越弱；}

        \textcolor{blue}{\cemh{[OH^-]} 越大，溶液碱性越强；\cemh{[OH^-]} 越小，溶液碱性越弱。}

        \textcolor{blue}{pH 越小，\cemh{[H+]} 越大，\cemh{[OH^-]} 越小，溶液酸性越强，碱性越弱；}

        \textcolor{blue}{pH 越大，\cemh{[H+]} 越小，\cemh{[OH^-]} 越大，溶液酸性越弱，碱性越强。}

        pH 的适用范围：\numrange{0}{14}

    \subsection{测量 pH 的方法}\label{subsec:测量pH的方法}
        \subsubsection{酸碱指示剂}\label{subsubsec:酸碱指示剂}

            指示剂的变色范围：

            \begingroup
            \setlength{\fboxsep}{0pt}
            % \small
            \newcommand{\boxstrut}{\rule[-0.2em]{0pt}{1.2em}}
            \begin{center}
            \begin{tabular}{lll}
                甲基橙 &\small 1\colorbox{red}{\makebox[4.2em]{\boxstrut 红色}}\colorbox{orange}{\makebox[2.6em]{\boxstrut 橙色}}\colorbox{yellow}{\makebox[15.2em]{\boxstrut 黄色}}12 & --3.1--4.4-- \\
                甲基红 &\small 1\colorbox{red}{\makebox[6.8em]{\boxstrut 红色}}\colorbox{orange}{\makebox[3.6em]{\boxstrut 橙色}}\colorbox{yellow}{\makebox[11.6em]{\boxstrut 黄色}}12 & --4.4--6.2-- \\
                石蕊   &\small 1\colorbox{red}{\makebox[7em]{\boxstrut 红色}}\colorbox{violet!85!white}{\makebox[7.6em]{\boxstrut \color{white}紫色}}\colorbox{blue}{\makebox[7.4em]{\boxstrut \color{white}蓝色}}12 & --4.5--8.3-- \\
                酚酞   &\small 1\colorbox{gray!25!white}{\makebox[14em]{\boxstrut 无色}}\colorbox{red!50!white}{\makebox[4em]{\boxstrut 粉红色}}\colorbox{red}{\makebox[4em]{\boxstrut 红色}}12 & --8.0--10.0-- \\
            \end{tabular}
            \end{center}
            \endgroup

        \subsubsection{pH 试纸}\label{subsubsec:pH试纸}

            广泛 pH 试纸（普通 pH 试纸）：只能测出\numrange{1}{14}范围的整数值

            精密 pH 试纸（有适用范围，读数精确到 \( 0.1 \) ）

            pH 试纸的使用注意事项——不可润湿

        \subsubsection{pH 计（酸度计）}\label{subsubsec:pH计}
            \textcolor{red}{测定前用蒸馏水洗净探头}

            直接读数，精确到 \( 0.01 \)

    \subsection{有关 pH 的计算}\label{subsec:有关pH的计算}

        \subsubsection{\texorpdfstring{\cemh{[H+]}}{H⁺}与 pH 之间的换算}\label{subsubsec:[H+]与pH之间的换算}
            \begin{xmp}
                [pH 计算例一]
                {}
                []
                []
                \vspace{-\baselineskip}
                \begin{enumerate}[itemindent=0em]
                    \item 已知\(c(\cemh{Cl-}) = \SI{0.01}{\mole\per\litre}\)，
                求其\(\cemh{pH}=\)\uline{\makebox[1.5em]{2}}，\(\cemh{[OH^-]}=\)\newline\uline{\makebox[9em]{\SI{1E-12}{\mole\per\litre}}}；
                    \item \(\cemh{pH} = 3\) 的稀硫酸的溶液浓度为\uline{\makebox[4.5em]{\SI{5E-4}{}}}\SI{}{\mole\per\litre}。
                \end{enumerate}
            \end{xmp}
            \begin{xmp}
                [pH 计算例二]
                {}
                []
                []
                \vspace{-\baselineskip}
                \begin{enumerate}
                    \item \SI{50}{\ml} \(\cemh{pH} = 1\)的盐酸溶液中，\( \conc[n]{H+} =\uline{\makebox[6em]{\SI{0.005}{\mole}}}\)；
                    \item 甲溶液的\(\cemh{pH} = 2\)，乙溶液的\(\cemh{pH} = 5\)，则溶液中 \cemh{[H+]} 甲是乙的\uline{\makebox[4em]{1000}}倍，\cemh{[OH^-]} 甲是乙的\uline{\makebox[4em]{\(\frac{1}{1000}\)}}倍。
                \end{enumerate}
            \end{xmp}
            
        \subsubsection{\texorpdfstring{\cemh{[OH^-]}}{OH⁻}与 pH 之间的换算}\label{subsubsec:[OH^-]与pH之间的换算}
            \begin{xmp}
                [pH 计算例三]
                {}
                []
                []
                \vspace{-\baselineskip}
                \begin{enumerate}
                    \item \SI{0.05}{\mole\per\litre}的氢氧化钡溶液，\(\cemh{pH}=\uline{\makebox[1.5em]{13}}\)；
                    \item 配置\(\cemh{pH} = 14\)的 \cemh{NaOH} 溶液\SI{100}{\ml}，需 \cemh{NaOH}\uline{\makebox[1.5em]{2}}\SI{}{\gram}；
                    \item \SI{5.6}{\gram}\,\cemh{KOH} 溶于水配成\SI{1}{\litre}溶液，\(\cemh{[OH^-]}=\uline{\makebox[5.5em]{\SI{0.1}{\mole\per\litre}}}\)，\(\cemh{pH}=\uline{\makebox[1.3em]{13}}\) 。
                \end{enumerate}
            \end{xmp}
            
        \subsubsection{稀释计算}\label{subsubsec:稀释计算}
            \begin{xmp}
                [pH 计算例四]
                {}
                []
                []
                取\SI{1}{\ml}\,\(\cemh{pH}=2\)的盐酸加水稀释到\SI{100}{\ml}后，该溶液的\(\cemh{pH}=\uline{\makebox[1.5em]{4}}\)；

                取\SI{1}{\ml}\,\(\cemh{pH}=12\)的 \cemh{NaOH} 溶液加水稀释到\SI{100}{\ml}后，该溶液的\(\cemh{pH}=\uline{\makebox[1.5em]{10}}\)。
            \end{xmp}
            \begin{xmp}
                [pH 计算例五]
                {}
                []
                []
                将\(\cemh{pH}=5\)的硫酸稀释\(10\)倍，\(\cemh{pH}=\uline{\makebox[1.5em]{6}}\)；

                将\(\cemh{pH}=5\)的硫酸稀释\(1000\)倍，\(\cemh{pH}=\uline{\makebox[1.5em]{7}}\)。
            \end{xmp}

            \textcolor{red}{[小结]：强酸强碱每稀释 \( 10 \) 倍，\cemh{[H+]}、\cemh{[OH^-]} 变化 \( 10 \) 倍，pH 改变 \( 1 \)；}

            \phantom{[小结]：}\textcolor{red}{但酸不会稀释成碱，碱不会稀释成酸，pH 最终逼近 \( 7 \)。}

        \subsubsection{混合计算}\label{subsubsec:混合计算}

        \noindent   1. 酸酸混合

            \stepcounter{example}%例六缺失
            \begin{xmp}
                [pH 计算例七]
                {}
                []
                []
                \(\cemh{pH}=2\)的盐酸\SI{100}{\ml}与 \(\cemh{pH}=3\)的硫酸\SI{200}{\ml}混合后，\(\cemh{pH}=\uline{\makebox[2em]{2.40}}\)。
            \end{xmp}

        \noindent   2. 碱碱混合

            \begin{xmp}
                [pH 计算例八]
                {}
                []
                []
                \(\cemh{pH}=10\)的 \cemh{Ba(OH)2} 溶液与 \(\cemh{pH}=12\)的 \cemh{NaOH} 溶液等体积混合后，溶液中的\(\cemh{[H+]}=\uline{\makebox[9.5em]{\SI{1.96E-12}{\mole\per\litre}}}\)，\(\cemh{pH}=\uline{\makebox[3em]{11.70}}\) 。
            \end{xmp}

        \noindent   3. 酸碱中和

            \begin{xmp}
                [pH 计算例九]
                {}
                []
                []
                \(\cemh{pH}=3\)的硫酸与 \(\cemh{pH}=10\)的 \cemh{NaOH} 溶液混合，若要混合后的溶液\(\cemh{pH}=7\)（恰好完全中和），则硫酸与 \cemh{NaOH} 溶液的体积比为\uline{\makebox[3em]{1:10}} 。
            \end{xmp}

            \begin{xmp}
                [pH 计算例十]
                {}
                []
                []
                把\SI{100.5}{\ml}\,\SI{0.2}{\mole\per\litre}的 \cemh{NaOH} 溶液加入到\SI{99.5}{\ml}\,\SI{0.1}{\mole\per\litre}的硫酸中，充分反应后所得溶液的\(\cemh{pH}=\uline{\makebox[2em]{11}} \)。
            \end{xmp}

        \noindent   4. 其它反应

            \begin{xmp}
                [pH 计算例十一]
                {}
                []
                []
                \SI{0.01}{\mole\per\litre}的硫酸\SI{5}{\ml}与\SI{0.01}{\mole\per\litre}的 \cemh{BaCl2} 溶液\SI{5}{\ml}混合，充分反应后所得溶液的\(\cemh{pH}=\uline{\makebox[2em]{2}}\)。
            \end{xmp}

            \textcolor{red}{[小结]：若反应不涉及 \cemh{H+}、\cemh{OH^-}，相当于稀释}

        \subsubsection{非常温计算}\label{subsubsec:非常温计算}

            \begin{xmp}
                [pH 计算例十二]
                {}
                []
                []
                某温度下，\cemh{NaCl} 溶液\(\cemh{pH}=6.5\)；则该温度下\(\cemh{pH}=2\)的硫酸与\(\cemh{pH}=\uline{\makebox[2.5em]{11}}\)的 \cemh{NaOH} 溶液等体积混合后恰好完全中和。该温度下\(\cemh{pH}=2\)的硫酸与\(\cemh{pH}=10\)的 \cemh{NaOH} 溶液按体积比\(\uline{\makebox[1.5em]{9}}:\uline{\makebox[3em]{101}}\)，混合后溶液\(\cemh{pH}=9\)。
            \end{xmp}

            \textcolor{red}{[小结]：先算出该温度下的\(K_\text{W}\)，其他计算同前。}
            
    \subsection{水电离出的\texorpdfstring{\cemh{[H+]}}{H⁺}和\texorpdfstring{\cemh{[OH^-]}}{OH⁻}\texorpdfstring{\footnotemark}{}}\label{subsec:水电离出的[H+]和[OH^-]}\footnotetext{暂时只有酸和碱}

        〔常温下〕某溶液中水电离出的\(\cemh{[H+]}=\SI{1E-11}{\mole\per\litre}\)，该溶液\(\cemh{pH}=\uline{\makebox[5em]{3 或 11}}\)；

        〔常温下〕某溶液中水电离出的\(\cemh{[OH^-]}=\SI{1E-11}{\mole\per\litre}\)，该溶液\(\cemh{pH}=\uline{\makebox[5em]{3 或 11}}\)。

        1、pH＝3 的盐酸溶液中，水电离出的\(\cemh{[H+]}_\text{水}=\uline{\makebox[9em]{\SI{1E-12}{\mole\per\litre}}}\)，水电离出的\(\cemh{[OH^-]}_\text{水}=\uline{\makebox[9em]{\SI{1E-12}{\mole\per\litre}}}\) 。

        2、pH＝11 的 \cemh{NaOH} 溶液中，水电离出的\(\cemh{[H+]}_\text{水}=\uline{\makebox[9em]{\SI{1E-12}{\mole\per\litre}}}\) ，水电离出的\(\cemh{[OH^-]}_\text{水}=\uline{\makebox[9em]{\SI{1E-12}{\mole\per\litre}}}\) 。

        
        \textcolor{red}{[小结]：酸碱抑制水的电离，使水的电离程度减小；}

        \textcolor{red}{若酸、碱电离出的 \cemh{[H+]} 和 \cemh{[OH^-]} 对应相等，则抑制程度相同。}

        \textcolor{red}{酸、碱溶液中，水电离的量看溶液的 \cemh{[H+]}、\cemh{[OH^-]} 中小的那个。}
